% FILE: 06_contribution_to_thesis.tex
% PROJECT: gaps — Structural Gap Registry and Closure Tracking

\section{Contribution to the Dissertation}
\label{sec:gaps-contribution}

\subsection{Contribution to Prediction vs.~Coordination}
\label{sec:gaps-contribution-prediction}

The gap registry makes a specific contribution to the dissertation's treatment of the
prediction-versus-coordination distinction. The old-regime process mining paradigm treats
conformance checking as a post-hoc prediction exercise: given an observed trace and a
reference model, predict whether the trace conforms. This prediction is made after the
fact, using approximation methods (token replay) that produce fitness values lacking
tight bounds.

The gap registry documents the gap between this prediction-oriented architecture and a
coordination-oriented one. In a coordination-oriented architecture, the type system
coordinates the admission of evidence before it reaches the execution engine. An
algorithm that accepts only \texttt{Evidence<EventLog, Admitted, W>} at its boundary
is not predicting whether the input is valid — it is enforcing, at compile time, that
the input has already been admitted under named laws. GAP\_001 documents the precise
engineering delta between these two paradigms: the absence of wasm4pm-compat in the
wasm4pm dependency graph. The transition from prediction to coordination is not
philosophical; it is a Cargo.toml edit and a set of function signature changes, with
a five-condition checklist in GAP\_001 specifying exactly what constitutes the
completed transition.

GAP\_004 extends this contribution: alignment-based conformance with
\texttt{Metric<FITNESS, N, D>} bounded by \texttt{Between01} replaces token-replay
approximation with exact optimal-alignment fitness. This is a shift from a predicted
approximation to a coordinated measurement with known bounds.

\subsection{Contribution to Process Evidence}
\label{sec:gaps-contribution-evidence}

The gap registry contributes to the dissertation's theory of process evidence by
identifying the precise locations where the evidence lifecycle is broken. Process
evidence is not merely data about a process — it is data with a documented chain of
custody, from raw observation through admission under named laws through execution in
a type-safe engine through conformance assessment with bounded metrics. The gap registry
maps every break in this chain:

\begin{itemize}
  \item GAP\_001 breaks the chain at the admission-to-execution boundary.
  \item GAP\_002 breaks the refusal surface — string-typed errors cannot be traced to
        named laws and cannot participate in evidence chains.
  \item GAP\_007 breaks the WF-net attestation link — a forgeable attestation path
        means WF-net soundness proofs are not trustworthy receipt artifacts.
  \item GAP\_008 breaks the compile-fail fixture chain — E0425 failures are import
        errors, not type-law evidence.
\end{itemize}

The contribution is not merely identifying these breaks. It is establishing the
evidentiary standard that determines what counts as a repair: the five-condition
GAP\_001 checklist, the named error code requirements of GAP\_008 (E0308, E0599, E0277,
E0080), and the non-forgeability requirement of GAP\_007 (\texttt{pub(crate)} or
\texttt{\#[deprecated]} gating of \texttt{attest\_witnessed()}).

\subsection{Contribution to Manufacturing Architecture}
\label{sec:gaps-contribution-manufacturing}

The gap registry contributes to the CodeManufactory manufacturing architecture by
demonstrating what a proof gate means in practice. A proof gate is not a runtime check
that can be bypassed by providing a conforming input — it is a type-system constraint
that cannot be bypassed at all without violating a structural law. The gap registry
documents the delta between the current state (runtime checks producing
\texttt{ValidationError(String)}) and the required state (compile-time type-law gates
producing \texttt{Refusal<R, W>} with named law type \texttt{R}).

The dependency graph in \texttt{GAP\_PRIORITY\_MATRIX.md} is itself a manufacturing
architecture artifact: it specifies the partial order over proof-gate manufacturing
steps, from GAP\_007 (lowest effort, independent) through GAP\_001 (critical-path
blocker) to the P2 and P3 gaps that can proceed in parallel once the bridge is in place.
The critical path \texttt{GAP\_007 $\rightarrow$ GAP\_001 $\rightarrow$ GAP\_002}
is the minimum manufacturing sequence required before any board-level process evidence
claim can be asserted with receipt-chain support.

\subsection{Contribution to Capital-Flow Infrastructure}
\label{sec:gaps-contribution-capital-flow}

\textbf{AUTHOR THESIS:} The capital-flow and settlement infrastructure described in the
thesis lineage (Post-Cyberpunk PCP, AI XYNZ, DAO settlement) requires that process
evidence flows through a verifiable chain from observation to settlement authorization.
A settlement authorization that cannot be traced to admitted, receipted process evidence
is not a valid settlement — it is an assertion without grounding.

The gap registry contributes to this layer by documenting precisely why the current
wasm4pm stack cannot support settlement-grade process evidence. GAP\_001's impact
assessment states that the type-law enforcement is zero, the admission gate is missing,
the witness discipline is missing, the lifecycle state is missing, and the receipt chain
is broken. A system with these defects cannot produce process evidence that meets the
evidentiary standard required for capital-flow authorization. The gap registry therefore
defines the engineering preconditions for capital-flow infrastructure deployment: all
eight named gaps must be formally closed, with receipts and resolution addendums, before
the wasm4pm stack can be claimed as settlement-grade process evidence infrastructure.
